\chapter{Web Application}
\section{Technology stack}

The application is built with:

\begin{itemize}
    \item \textbf{Backend:} Python 3 + Flask
    \item \textbf{Database Driver:} oracledb
    \item \textbf{Frontend:} HTML5, CSS3, Jinja2 templates
    \item \textbf{Database:} Oracle Object-Relational DB
\end{itemize}

\section{Project structure}

\begin{verbatim}
webapp/
|-- app.py
|-- database.py
|-- static/
|   |-- css/
|   |   `-- style.css
`-- templates/
    |-- base.html
    |-- index.html
    |-- register_customer.html
    |-- add_booking.html
    |-- add_event_location.html
    |-- view_location_team.html
    `-- ranked_locations.html
\end{verbatim}

\section{Local execution}

The Flask application can be started after preparing the Oracle schema with
\texttt{script.sql} and test data.

\begin{lstlisting}[language=bash]
cd webapp
pip install flask oracledb
set DB_USER=TEST
set DB_PASSWORD=test123
set DB_DSN=localhost:1521/XE
python app.py
\end{lstlisting}

Then open \texttt{http://127.0.0.1:5000} in the browser.

\section{Database connection layer}

The \texttt{database.py} module centralizes Oracle connectivity:

\begin{lstlisting}[language=Python]
import os
import oracledb

def get_connection():
    user = os.getenv('DB_USER', 'stageup_user')
    password = os.getenv('DB_PASSWORD', 'stageup_pwd')
    dsn = os.getenv('DB_DSN', 'localhost:1521/XEPDB1')
    return oracledb.connect(user=user, password=password, dsn=dsn)
\end{lstlisting}

\section{Route mapping}

\begin{table}[H]
    \centering
    \resizebox{\textwidth}{!}{
    \begin{tabular}{|c|p{3.0cm}|c|p{3.5cm}|p{3.5cm}|}
    \hline
    \textbf{Op.} & \textbf{Endpoint} & \textbf{Method} & \textbf{DB object} & \textbf{Template} \\
    \hline
    1 & \texttt{/register\_customer} & GET, POST & \texttt{proc\_register\_customer} & \texttt{register\_customer.html} \\
    \hline
    2 & \texttt{/add\_booking} & GET, POST & \texttt{proc\_add\_booking} & \texttt{add\_booking.html} \\
    \hline
    3 & \texttt{/add\_event\_location} & GET, POST & \texttt{proc\_add\_event\_location} & \texttt{add\_event\_location.html} \\
    \hline
    4 & \texttt{/view\_location\_team} & GET, POST & \texttt{func\_get\_team\_by\_location} & \texttt{view\_location\_team.html} \\
    \hline
    5 & \texttt{/ranked\_locations} & GET & \texttt{vw\_ranked\_locations} & \texttt{ranked\_locations.html} \\
    \hline
    \end{tabular}
    }
    \caption{Route mapping summary}
\end{table}

\noindent All write routes (Operations 1--3) follow the same flow: read form data, call PL/SQL object, commit, flash message, and redirect.

\section{UI screenshots}

\begin{figure}[H]
    \centering
    \safeincludegraphics[width=0.95\textwidth]{images/home.png}
    \caption{StageUp Web Application - Homepage}
\end{figure}

\begin{figure}[H]
    \centering
    \begin{subfigure}[t]{\textwidth}
        \centering
        \safeincludegraphics[width=0.95\textwidth]{images/op1app.png}
        \caption{Operation 1 - Form entry}
    \end{subfigure}
    \vspace{0.5cm}
    \begin{subfigure}[t]{\textwidth}
        \centering
        \safeincludegraphics[width=0.95\textwidth]{images/op1apppostins.png}
        \caption{Operation 1 - After submission}
    \end{subfigure}
    \caption{Operation 1 - Register Customer}
\end{figure}

\begin{figure}[H]
    \centering
    \begin{subfigure}[t]{\textwidth}
        \centering
        \safeincludegraphics[width=0.95\textwidth]{images/op2app.png}
        \caption{Operation 2 - Form entry}
    \end{subfigure}
    \vspace{0.5cm}
    \begin{subfigure}[t]{\textwidth}
        \centering
        \safeincludegraphics[width=0.95\textwidth]{images/op2apppostins.png}
        \caption{Operation 2 - After submission}
    \end{subfigure}
    \caption{Operation 2 - Add Booking}
\end{figure}

\begin{figure}[H]
    \centering
    \begin{subfigure}[t]{\textwidth}
        \centering
        \safeincludegraphics[width=0.95\textwidth]{images/op3app.png}
        \caption{Operation 3 - Form entry}
    \end{subfigure}
    \vspace{0.5cm}
    \begin{subfigure}[t]{\textwidth}
        \centering
        \safeincludegraphics[width=0.95\textwidth]{images/op3apppostins.png}
        \caption{Operation 3 - After submission}
    \end{subfigure}
    \caption{Operation 3 - Add Event Location}
\end{figure}

\begin{figure}[H]
    \centering
    \begin{subfigure}[t]{\textwidth}
        \centering
        \safeincludegraphics[width=0.95\textwidth]{images/op4app.png}
        \caption{Operation 4 - Form entry}
    \end{subfigure}
    \vspace{0.5cm}
    \begin{subfigure}[t]{\textwidth}
        \centering
        \safeincludegraphics[width=0.95\textwidth]{images/op4apppostins.png}
        \caption{Operation 4 - After search}
    \end{subfigure}
    \caption{Operation 4 - View Team by Location}
\end{figure}

\begin{figure}[H]
    \centering
    \safeincludegraphics[width=0.95\textwidth]{images/op5app.png}
    \caption{Operation 5 - Ranked Locations (read-only view)}
\end{figure}

\section{Final notes}

\begin{itemize}
    \item Each exam operation has a dedicated route and template.
    \item Database logic is encapsulated in stored procedures/functions.
    \item The routes use \texttt{try/except/finally} blocks to guarantee cursor/connection cleanup.
    \item Operation 4 returns either the team name or an error message generated by the database layer.
    \item Credentials are loaded from environment variables in \texttt{database.py}.
\end{itemize}







